\section{Physical consequences of the square-grid correspondence}

\begin{theorem}[Equality of physical vectors and errors]
    \label{thm:peps_square_vector_error}
    \lean{TNLean.PEPS.Approximation.squareGridStateIsometry,TNLean.PEPS.Approximation.vector_pinnedTensorToGraphTensor,TNLean.PEPS.Approximation.vector_vectorTensorToGraphTensor,TNLean.PEPS.Approximation.norm_pinnedTensorToGraphTensor,TNLean.PEPS.Approximation.nonzero_pinnedTensorToGraphTensor,TNLean.PEPS.Approximation.normalized_error_pinnedTensorToGraphTensor,TNLean.PEPS.Approximation.normalized_error_vectorTensorToGraphTensor}
    \leanok
    \uses{thm:peps_square_grid_contraction}
    The physical coordinates are identical in both presentations, so the coordinate
    identification is a linear isometry. The transported contraction equals the
    original vector $\Phi$. In particular its norm, nonvanishing, and every error
    $\|\|\Phi\|^{-1}\Phi-z\Omega\|$ are unchanged, for every $z\in\C$.
    The equality includes $\Phi=0$ with the convention $0^{-1}=0$.
\end{theorem}
\begin{proof}
    \leanok
    \uses{thm:peps_square_grid_contraction}
    Apply coordinatewise equality of the contractions in their common Euclidean space.
\end{proof}

\begin{theorem}[Dimension constraints and crossing edges]
    \label{thm:peps_square_bonds_boundary}
    \lean{TNLean.PEPS.Approximation.graphBondDim_forall,TNLean.PEPS.Approximation.vectorTensorToGraphTensor_bondDim_pos,TNLean.PEPS.Approximation.maxBondDim_le_iff,TNLean.PEPS.Approximation.forwardSquareBoundaryEquiv,TNLean.PEPS.Approximation.forwardSquareBoundary_card}
    \leanok
    \uses{def:peps_square_edge_correspondence}
    Any pointwise condition on $D_e$ is equivalent to the same condition on the
    transported dimensions. For $B\in\N$, the maximum bond dimension is at most
    $B$ exactly when each transported bond dimension is at most $B$.
    For every region $R\subseteq\Lambda$, the edge correspondence restricts to
    a bijection of crossing edges and preserves their cardinality.
\end{theorem}
\begin{proof}
    \leanok
    \uses{def:peps_square_edge_correspondence}
    Every edge retains its dimension and its two endpoints. The empty maximum is zero.
\end{proof}

\begin{theorem}[Regional density and entropy]
    \label{thm:peps_square_regional_density}
    \lean{TNLean.PEPS.Approximation.squareComplementConfigEquiv,TNLean.PEPS.Approximation.regionReducedDensity_eq_piSubtype,TNLean.PEPS.Approximation.regionReducedDensity_eq_pinnedReducedDensity,TNLean.PEPS.Approximation.regionReducedDensity_pinnedTensorToGraphTensor,TNLean.PEPS.Approximation.regionReducedDensity_vectorTensorToGraphTensor,TNLean.PEPS.Approximation.entropy_pinnedTensorToGraphTensor,TNLean.PEPS.Approximation.entropy_vectorTensorToGraphTensor}
    \leanok
    \uses{thm:peps_square_grid_contraction}
    For every region $R$, the two contractions have equal unnormalized reduced
    densities $\tr_{R^c}|\Phi\rangle\langle\Phi|$ and equal von Neumann entropies.
    Writing the complement as a finite-set difference or as nonmembership gives
    the same partial trace after the coordinatewise identification of complement
    configurations.
\end{theorem}
\begin{proof}
    \leanok
    \uses{thm:peps_square_grid_contraction}
    Substitute the coefficient equality into the outer product and its partial trace.
    The change of complement coordinates is a bijective reindexing of the finite sum.
    Apply congruence of von Neumann entropy in its matrix argument.
\end{proof}

\begin{definition}[The source approximation constraints]
    \label{def:peps_source_approximation_constraints}
    \lean{TNLean.PEPS.Approximation.Pinned.HasPEPSApproximation,TNLean.PEPS.Approximation.Vector.PhaseErrorAtMost}
    \leanok
    \uses{def:peps_square_local_presentations}
    The physical-first approximation constraint asks for positive $D_e\le CL^c$,
    a nonzero contraction $\Phi$, and a phase $\theta$ with normalized vector error
    at most $L^{-1}$. The virtual-first phase-error constraint additionally requires
    that the chosen phase minimize the error over all real phases.
\end{definition}

\begin{theorem}[Transfer of supplied approximation witnesses]
    \label{thm:peps_square_approximation_transfer}
    \lean{TNLean.PEPS.Approximation.phaseErrorAtMost_vectorTensorToGraphTensor,TNLean.PEPS.Approximation.hasPEPSApproximation_of_pinned,TNLean.PEPS.Approximation.hasPEPSApproximation_of_vector}
    \leanok
    \uses{def:peps_source_approximation_constraints,thm:peps_square_vector_error,thm:peps_square_bonds_boundary}
    An approximation witness in either source presentation gives a graph PEPS on
    the original open square with the same positive dimensions, bound $CL^c$,
    and error $L^{-1}$. The entire minimizing-phase condition is preserved for the
    virtual-first presentation.
\end{theorem}
\begin{proof}
    \leanok
    \uses{thm:peps_square_vector_error,thm:peps_square_bonds_boundary}
    Use the transported tensor. Its dimensions are identical edge by edge; substitute
    the equality of its vector in the nonzero and phase-error conditions.
\end{proof}


\begin{definition}[Physical-first regional coefficient matrix]
    \label{def:peps_pinned_regional_matrix}
    \lean{TNLean.PEPS.Approximation.Pinned.RegionConfiguration,TNLean.PEPS.Approximation.Pinned.joinConfigurations,TNLean.PEPS.Approximation.Pinned.coefficientMatrix,TNLean.PEPS.Approximation.Pinned.reducedDensity}
    \leanok
    For configurations $x$ on $R$ and $z$ on $R^c$, let $x\cup z$ be their
    joined configuration and set $M_{x,z}=\Phi(x\cup z)$. The unnormalized regional
    density is $MM^\dagger$, which equals the partial trace in
    Theorem~\ref{thm:peps_square_regional_density}.
\end{definition}
